\documentclass[aspectratio=169,UTF8]{ctexbeamer}
\usetheme{Madrid}
\usecolortheme{default}
\definecolor{DeepBlue}{HTML}{2563EB}
\definecolor{DeepInk}{HTML}{172033}
\definecolor{DeepGreen}{HTML}{059669}
\setbeamercolor{structure}{fg=DeepBlue}
\setbeamercolor{title}{fg=white,bg=DeepBlue}
\setbeamercolor{frametitle}{fg=white,bg=DeepBlue}
\setbeamercolor{alerted text}{fg=DeepGreen}
\setbeamertemplate{navigation symbols}{}
\setbeamertemplate{footline}[frame number]
\usepackage{booktabs}
\usepackage{listings}
\usepackage{xcolor}
\usepackage{hyperref}
\hypersetup{colorlinks=true,urlcolor=DeepBlue}
\lstset{
  basicstyle=\ttfamily\scriptsize,
  keywordstyle=\color{DeepBlue},
  stringstyle=\color{DeepGreen},
  commentstyle=\color{gray},
  showstringspaces=false,
  breaklines=true,
  frame=single,
  columns=fullflexible
}
\newcommand{\api}[1]{\texttt{#1}}
\newcommand{\keypoint}[1]{\alert{#1}}
\AtBeginSection[]{
  \begin{frame}{本节路线}
    \tableofcontents[currentsection]
  \end{frame}
}

\title{第11课：Anthropic Messages API}
\subtitle{生态迁移、内容块、工具与流式用量}
\author{DeepSeek API 全栈课程}
\date{文档快照：2026-11-05}
\begin{document}
\frame{\titlepage}
\begin{frame}{学习目标}
\begin{itemize}
  \item 用 Anthropic SDK 调用 DeepSeek
  \item 理解模型映射与字段兼容性
  \item 实现 tool\_use / tool\_result 循环
  \item 正确读取流式事件、缓存与服务端工具用量
\end{itemize}
\end{frame}
\section{接入基础}
\begin{frame}{地址与鉴权}
base URL：\api{https://api.deepseek.com/anthropic}。可用 \api{x-api-key} 或 \api{Authorization: Bearer}。\\[8pt]
Anthropic SDK 要求显式 max\_tokens；裸 HTTP 可省略，由模型/模式决定默认值。
\end{frame}
\begin{frame}{模型名映射}
\begin{itemize}
  \item claude-opus* $\rightarrow$ deepseek-v4-pro（按 Pro 计费）
  \item claude-haiku* / claude-sonnet* $\rightarrow$ deepseek-flash
  \item 直接使用 deepseek-flash / deepseek-v4-pro 最清晰
  \item 非 claude 前缀的未知模型名返回 400
\end{itemize}
映射用于兼容受模型名限制的客户端，不应掩盖真实成本。
\end{frame}
\begin{frame}[fragile]{最小调用}
\begin{lstlisting}[language=Python]
client = anthropic.Anthropic(
  api_key=os.environ["DEEPSEEK_API_KEY"],
  base_url="https://api.deepseek.com/anthropic",
)
message = client.messages.create(
  model="deepseek-flash", max_tokens=500,
  system="回答简洁。",
  messages=[{"role":"user", "content":"你好"}],
)
\end{lstlisting}
\end{frame}
\section{字段兼容}
\begin{frame}{常用字段}
system、messages、stop\_sequences、stream 完全支持；metadata 仅 user\_id 生效。thinking 支持 enabled/disabled，adaptive 视为 enabled，budget\_tokens 忽略；output\_config 仅 effort 生效。
\end{frame}
\begin{frame}{会被忽略的字段}
container、mcp\_servers、service\_tier、top\_k、cache\_control、text citations 等被忽略。思考模式 temperature 固定 1；top\_p 行为仍依 Flash/Pro。\\[8pt]
兼容成功不等于语义完全相同，应建立迁移测试。
\end{frame}
\begin{frame}{消息内容块}
支持 text、thinking、tool\_use、tool\_result、server\_tool\_use、web\_search\_tool\_result；不支持 document、search\_result、redacted\_thinking、MCP/代码执行结果等。\\[8pt]
Flash 支持 user 和 tool\_result 中的 image；Pro 替换为 \api{[Unsupported Image]}。
\end{frame}
\section{工具与思考}
\begin{frame}{客户端工具定义}
工具直接提供 name、description、input\_schema，\keypoint{不要给客户端函数传 type=custom}，否则 422。\\[8pt]
tool\_choice：none/auto 支持；any 按 auto；指定 tool 仅非思考模式支持。
\end{frame}
\begin{frame}{回传工具结果}
assistant 若含 tool\_use，紧接的消息必须是 user，并为所有调用提供匹配 tool\_use\_id 的 tool\_result；该 user 消息中所有 tool\_result 必须放在文本之前。
\end{frame}
\begin{frame}{思考工具链}
建议原样回传 assistant 完整 content 数组（thinking + tool\_use）。服务端有时能按调用 ID 恢复遗漏推理，但无法恢复会 400，\keypoint{不要依赖自动恢复}。
\end{frame}
\section{服务端工具与流}
\begin{frame}{Anthropic 服务端 web search}
声明 web\_search 服务端工具后，由服务端执行并继续回答，客户端无需 tool\_result。输出含 server\_tool\_use 与 web\_search\_tool\_result。\\[8pt]
\api{usage.server\_tool\_use} 区分尝试次数 web\_search\_requests 与成功有结果次数 tool\_search。
\end{frame}
\begin{frame}{流式事件}
message\_start $\rightarrow$ content\_block\_start/delta/stop $\rightarrow$ message\_delta $\rightarrow$ message\_stop，期间可有 ping 和保活注释；没有 \api{data: [DONE]}。\\[8pt]
最终 Token 以 message\_delta.usage 为准，不与 message\_start 相加。
\end{frame}
\begin{frame}{Anthropic usage}
总输入 = input\_tokens + cache\_read\_input\_tokens + cache\_creation\_input\_tokens；当前 cache creation 常为 0。流式中断且未收到 message\_stop，响应不完整。
\end{frame}
\begin{frame}{图片与 Files}
image.source.type 可为 base64、url、file。file 图片和 Anthropic Files API 必须带 \api{anthropic-beta: files-api-2025-04-14}；Files 路径、分页和字段差异见第8课。
\end{frame}
\begin{frame}{本课例程}
\begin{enumerate}
  \item \api{01\_anthropic\_basic.py}：基础消息与思考控制
  \item \api{02\_anthropic\_stream.py}：流式文本与最终用量
  \item \api{03\_anthropic\_tool.py}：工具内容块循环
\end{enumerate}
\end{frame}
\begin{frame}{官方资料}
\small
\url{https://api-docs.deepseek.com/zh-cn/guides/anthropic_api}\\
\url{https://api-docs.deepseek.com/zh-cn/api/create-message}
\end{frame}
\end{document}
